% Chương 1

\chapter{TỔNG QUAN ĐỀ TÀI}

\label{Chapter1}

%----------------------------------------------------------------------------------------

\newcommand{\keyword}[1]{\textbf{#1}}
\newcommand{\tabhead}[1]{\textbf{#1}}
\newcommand{\code}[1]{\texttt{#1}}
\newcommand{\file}[1]{\texttt{\bfseries#1}}
\newcommand{\option}[1]{\texttt{\itshape#1}}

%----------------------------------------------------------------------------------------

\section{Mở đầu}

Trong bức tranh tổng thể của nền y học hiện đại, chẩn đoán hình ảnh kỹ thuật số định hình như một công cụ then chốt, quyết định tính chính xác của quy trình phát hiện bệnh lý, theo dõi tiến triển lâm sàng và xây dựng phác đồ điều trị. Sự tiến bộ mang tính bước ngoặt của các công nghệ hình ảnh tiên tiến như chụp cắt lớp vi tính (CT), cộng hưởng từ (MRI), hay chụp cắt lớp phát xạ positron (PET) đã tái định nghĩa khả năng quan sát y khoa. Đặc biệt, sự ra đời của công nghệ chụp cắt lớp quang học và mạch máu quang học (\textit{Optical Coherence Tomography/Angiography} - OCT/OCTA) đã cho phép các chuyên gia y tế tiếp cận cấu trúc vi mạch và mô sinh học dưới dạng không gian ba chiều (3D) ở độ phân giải siêu cao. Khả năng hiển thị trực quan và chi tiết các khối bất thường này đóng vai trò tối quan trọng trong việc nâng cao chất lượng và tính kịp thời cho các quyết định lâm sàng.

Tuy nhiên, việc nâng cao độ phân giải hình ảnh lại đặt ra một thách thức lớn về mặt kỹ thuật do sự bùng nổ của khối lượng dữ liệu thô. Một khối dữ liệu ảnh 3D (Volume dữ liệu) biểu diễn hệ thống vi mạch thường bao phủ một không gian quy mô lớn, chứa từ hàng trăm triệu cho đến hơn một tỷ điểm dữ liệu (\textit{voxel}). Khi đối mặt với khối lượng thông tin đồ sộ này, việc triển khai chuỗi toán tử xử lý từ lọc nhiễu, phân mảnh cấu trúc cho đến trích xuất đặc trưng trên kiến trúc tính toán tuần tự của các bộ vi xử lý trung tâm (CPU) truyền thống lập tức bộc lộ những hạn chế cố hữu. Sự chậm trễ trong thời gian xử lý tạo ra một nút thắt cổ chai hệ thống nghiêm trọng, hoàn toàn không tương thích với yêu cầu chẩn đoán thời gian thực hoặc phản ứng nhanh trong các tình huống y tế khẩn cấp, đồng thời dẫn đến hiệu suất khai thác tài nguyên phần cứng thấp.

Để vượt qua rào cản hiệu năng của kiến trúc cũ, xu hướng dịch chuyển sang mô hình tính toán song song khối lượng lớn thông qua nền tảng công nghệ CUDA của NVIDIA trên các bộ xử lý đồ họa (GPU) được xem là một bước đi đột phá và mang tính tất yếu. Nhờ sở hữu kiến trúc phần cứng tối ưu cho phép điều phối và thực thi đồng thời hàng chục ngàn luồng dữ liệu (\textit{threads}), GPU mở ra khả năng tăng tốc vượt trội cho các tác vụ lặp đi lặp lại trên lưới ma trận đa chiều, từ đó thiết lập một đường ống xử lý dữ liệu (\textit{Data Pipeline}) liền mạch, đồng bộ cho các hệ thống chẩn đoán thông minh.

Hướng tới mục tiêu giải quyết triệt để bài toán hiệu năng tính toán trong phân tích ảnh y khoa, đề tài “Ứng dụng lập trình song song CUDA C++ trong xử lý và phân tích khối dữ liệu ảnh y khoa 3D” đã được nghiên cứu và thực hiện. Trên cơ sở thực nghiệm với khối ảnh mô phỏng từ bộ dữ liệu vi mạch mu bàn tay (\textit{Dorsum Hand OCT/OCTA}), nghiên cứu tập trung xây dựng một hệ thống xử lý tốc độ cao được tự động hóa hoàn toàn thông qua cấu trúc mã nguồn CUDA C++. Tiến trình xử lý của hệ thống bắt đầu bằng khâu giảm chiều không gian thông qua kỹ thuật ép khối dữ liệu ảnh 3D thành bản đồ bề mặt 2D (\textit{En-face/Maximum Intensity Projection}) để cung cấp góc nhìn tổng quan về mạng lưới mạch máu. Tiếp theo, hệ thống áp dụng các bộ lọc không gian cửa sổ trượt (Mean Blur) chạy song song song hành cùng thuật toán phân ngưỡng nhị phân nhằm triệt tiêu nhiễu nền và phân tách chính xác cấu trúc hình học của hệ thống vi mạch. Từ kết quả ảnh đã tiền xử lý, nhánh phân tích hình thái sẽ tự động đo lường mật độ mạng lưới mạch máu (\textit{Vessel Density}) thông qua các toán tử đếm nguyên tử (\textit{Atomic Operations}), trong khi nhánh phân tích kết cấu chịu trách nhiệm khởi tạo ma trận đồng xuất hiện mức xám (GLCM) kích thước $256 \times 256$ ngay trên bộ nhớ GPU để trích xuất các đặc trưng toán học chuyên sâu của mô sinh học bao gồm độ tương phản (Contrast), độ đồng nhất (Homogeneity) và năng lượng (Energy).
%----------------------------------------------------------------------------------------
